% Generated draft: requires mathematical and source review.
\documentclass[11pt]{article}
\usepackage[margin=25mm]{geometry}
\usepackage{fontspec}
\setmainfont{Latin Modern Roman}
\usepackage{amsmath,amssymb,mathtools,unicode-math}
\setmathfont{Latin Modern Math}
\usepackage{enumitem,needspace,xurl,hyperref}
\setlength{\parindent}{0pt}
\setlength{\parskip}{5pt}
\setlength{\emergencystretch}{3em}
\newcommand{\N}{\mathbb{N}}
\newcommand{\Z}{\mathbb{Z}}
\newcommand{\Q}{\mathbb{Q}}
\newcommand{\R}{\mathbb{R}}
\newcommand{\C}{\mathbb{C}}
\newcommand{\F}{\mathbb{F}}
\newcommand{\A}{\mathbb{A}}
\newcommand{\PP}{\mathbb{P}}
\newcommand{\E}{\mathbb{E}}
\newcommand{\eps}{\varepsilon}
\newcommand{\dd}{\,\mathrm{d}}
\begin{document}
\section*{UnsolvedMath ID: 30004495}
\textbf{Source problem code:} catalogue entry
\subsection{Definitions and mathematical statement}
Write \(\Phi\) for the period map mentioned in the source and \(L=K_B+[Z]\) for the indicated log canonical bundle, using additive notation for line bundles or divisor classes. The record does not define \(B\), \(Z\), the period map, or their required hypotheses; consequently, the following implications are statement schemata whose intended domain remains unresolved:
\[
\begin{aligned}
\Phi\text{ satisfies generic local Torelli}
&\ \Longrightarrow\ L\text{ is nef and big},\\
\Phi\text{ satisfies local Torelli}
&\ \Longrightarrow\ L\text{ is globally generated}.
\end{aligned}
\]
Both implications are asked to hold for every object satisfying the intended, but unspecified, ambient hypotheses.

For clarity, when \(B\) is a smooth projective variety and \(Z\) is a Cartier divisor, the conventional interpretation of this notation is
\[
L=\omega_B\otimes\mathcal O_B(Z).
\]
This describes a possible interpretation of the notation, not additional hypotheses supplied by the record. For a line bundle \(L\) on a projective variety \(B\) of dimension \(d\), nefness means \(\deg(L|_C)\geq 0\) for every integral projective curve \(C\subseteq B\), and bigness means
\[
\limsup_{m\to\infty}\frac{h^0(B,L^{\otimes m})}{m^d}>0.
\]
Here ``free'' is interpreted as globally generated: the evaluation map
\[
H^0(B,L)\otimes\mathcal O_B\longrightarrow L
\]
is surjective. This is stronger than semiampleness, which requires global generation only for some positive tensor power of \(L\).

For a holomorphic map between complex manifolds, the usual differential interpretation of local Torelli is injectivity of \(\mathrm d\Phi_b\) at every point, and generic local Torelli means such injectivity on a dense open subset. The supplied record does not specify whether these interpretations apply directly to its period map, including at the boundary, or how its target is to be understood.

The supplied dated assessment reports a result under a stronger ``local Torelli condition on the triple'': the log canonical bundle is nef and big, and some positive multiple is free. The triple and its stronger condition are not defined in the record. This reported result therefore cannot be identified here with either requested implication, and freeness of a positive multiple does not itself give freeness of \(L\).
\subsection{Short English statement}
The source asks whether generic local Torelli forces the log canonical bundle to be nef and big, and whether local Torelli forces the bundle itself to be globally generated. The ambient geometric hypotheses and the precise Torelli conditions are missing. The supplied partially solved label and stronger-condition result are dated curation requiring independent review.
\subsection{Sources}
\begin{itemize}
\item Supplied local UnsolvedMath registry record, permanent ID \(30004495\), source problem code \texttt{catalogue entry}, title: ``Positivity of Log Canonical Bundles under Local Torelli.'' The record belongs to the Oberwolfach Reports Open Problems collection and supplies proposed year 2020; no individual proposer is identified.
\item Supplied source citation: \emph{Komplexe Analysis -- Algebraicity and Transcendence} (2021), Oberwolfach Reports, pp.\ 1263--1313, DOI 10.4171/owr/2020/24. The literature-source label is ``Oberwolfach Report 24/2020.'' \url{https://doi.org/10.4171/owr/2020/24}.
\item Supplied literature source: \emph{Natural line bundles on completions of period mappings}. \url{https://arxiv.org/abs/2102.06310}. The supplied record gives no authors or theorem locator.
\item Literature assessment checked at \texttt{2026-08-22}: supplied dated curation, not independently verified. It labels the problem partially solved and reports nefness, bigness, and freeness of a positive multiple under a stronger local Torelli condition on the triple, while stating that the requested generic-local-Torelli and local-Torelli conjectures are not established in full generality. The background explicitly requests independent verification of time-sensitive status and cited claims.
\end{itemize}
\end{document}
